% problem-000238 9 芳香亲电取代反应
\section*{9. 芳香亲电取代反应}
\textit{下列为分问。}

\subsection*{9-1}
在芳香亲电取代反应中, 芳环上的取代基可分为致活基和致钝基; 或依据产物结构, 将取代基分为邻、对位定位基和间位定位基。确定以下哪些基团既是致钝基团, 又是邻对位定位基团。(   )

(a) Cl; (b) Me; (c) NHCOMe; (d) CONHMe; (e) Br; (f) 以上基团都不是。

\subsection*{9-2}
在某些情况下，烷基作为邻对位定位基团也无法决定产物的最终区域选择性。如，苯的 Friedel-Craft 乙基化反应生成多种三乙基苯的混合物，其中 1,3,5-三乙基苯产率随着反应进行会发生较显著的变化。

\subsection*{9-2}
-1 随反应时间延长，1,3,5-三乙基苯的产率将增多还是减少？

\subsection*{9-2}
-2 1,3,5-三乙基苯的产率随反应时间改变的原因是什么？

\subsection*{9-3}
致钝基团对反应速率的影响并不完全取决于其吸电子效应。如, 3-氟-1,2,4,5-四甲基苯在 $30^{\circ} \mathrm{C}$ 乙酸溶液中发生溴代反应的速率是 1,2,4,5-四甲基苯的 2.31 倍。利用相关中间体的结构说明为何 3-氟-1,2,4,5-四甲基苯反应速率更快。

\subsection*{9-4}
以下三个化合物在 $\mathrm{NH_4NO_3}$ /TFAA/MeCN条件下进行单硝化反应:

（图：）
1

（图：）

2

（图：）

3

依据以上信息，回答以下问题：

\subsection*{9-4}
-1 在这三个化合物中，哪个化合物的硝化反应速率最快：（ ）。

\subsection*{9-4}
- 2 在反应最快的化合物中，以下说法哪一个是正确的：（）

(a) A 环的反应速率快； (b) B 环的反应速率快； (c) A 环和 B 环的反应速率相同。

\subsection*{9-4}
-3 画出反应最快的化合物进行硝化反应的关键中间体。

\subsection*{9-4}
-4 依据以上信息，完成以下反应式：

\subsection*{9-4}
-4-1

（图：）

\subsection*{9-4}
-4-2

（图：）

\subsection*{9-5}
文献报道了以下 Friedel-Craft 酰基化反应:

4

（图：）

\subsection*{9-5}
-1 依据以上信息完成以下反应式 (提示: 只需要画出一个主产物芳基酮):

（图：）

\subsection*{9-5}
-2 化合物 5 和 7 均能得到两个电子，分别画出这两个芳基酮得到两个电子后的结构式，并判断哪一个更稳定。
